\documentclass[]{article}
% {---
\usepackage{multicol}
\usepackage[T1]{fontenc}
\usepackage[margin=1.00in]{geometry}

\renewcommand\familydefault{\sfdefault}

\usepackage{graphicx}
\graphicspath{
    {/Users/laurel/Dropbox/Figures/}
    {/Users/laurel/Dropbox/Figures/Lightcurves/}
    {/Users/laurel/Dropbox/Figures/Temp/}
}

\setlength{\parindent}{0em} % default is 15 pt.
\setlength{\parskip}{6pt}

% Headers & Footers
\usepackage{fancyhdr}
\setlength{\headheight}{15.0pt}
\pagestyle{fancy}
\fancyhf{}
\fancyhead[R]{\footnotesize\leftmark}
\renewcommand{\headrulewidth}{0.5pt} % remove the header rule
\renewcommand*{\headrule}{\color{black!50!white}}
\cfoot{\thepage}
\pagenumbering{gobble}

\usepackage{amsmath}
%\usepackage{mathtools}
%\usepackage{mathptmx}  % times font (including math)
%\usepackage{mathpazo}  % palatino font (including math)
\usepackage{xcolor}
\usepackage{ragged2e}
% ---}

% Section titles/labels
% {---
\usepackage{titlesec}
%\titleformat{<command>}
%   [<shape>]{<format>}{<label>}{<sep>}{<before-code>}[<after-code>]
\titleformat{\section}
    {\fontsize{16}{14}\selectfont\centering\bfseries\scshape}{\thesection.}{1em}{}
\titleformat{\subsection}
    {\fontsize{13}{15}\selectfont\centering\bfseries}{\thesubsection.}{1em}{}
\titleformat{\subsubsection}
    {\centering\bfseries\itshape}{\thesubsubsection.}{1em}{}
\titleformat{\paragraph}
    {\bfseries\itshape}
    {}
    {}
    {}
    [\vspace{-1.0em}]

%\newcounter{paranum} % Counter for paragraphs
%\titleformat{\paragraph}
%    {\centering\itshape}
%    {\thesubsubsection.\refstepcounter{paranum}\theparanum.}{1em}{\vspace{-1.0em}}
%\titleformat{\subparagraph}
%    {\centering\itshape}
%    {}{1em}{\vspace{-1.0em}}
\setcounter{secnumdepth}{3}
% ---}

% List environments
% {---
\usepackage{enumitem}
\setlist[]{noitemsep}
% ---}

% Margin notes
% {---

\usepackage{tikz}
\usepackage[color=white]{background}  % No background at first...
\usepackage{schemata}  % braces around text in margin
\usepackage{marginnote}
%\renewcommand*{\marginfont}{\small\sffamily}

\definecolor{bred}{rgb}{0.8, 0.0, 0.0}
\definecolor{purple}{rgb}{0.44, 0.16, 0.39}
\definecolor{blue}{rgb}{0.0, 0.2, 0.6}
\definecolor{green}{rgb}{0.0, 0.42, 0.24}

% Action required vs. simple reference note?
\newcommand{\mynote}[2]{\hspace{0pt}\reversemarginpar\marginnote{\raggedright
    %\textcolor{green}{\Large\hfill$\star\;\longrightarrow$\newline}
    {\footnotesize#1}\newline
    \textcolor{black}{\scriptsize #2}
%\newcommand{\mynote}[2]{\hspace{0pt}\reversemarginpar\marginnote{\raggedright
%    \textcolor{green}{\footnotesize$\star$ #1\newline}
%    \textcolor{black}{\scriptsize #2}
    %\textcolor{green}{\fontsize{9}{10.0}\selectfont$\leadsto$ #1}
    %\textcolor{black}{\fontsize{7}{8.0}\selectfont #2}
}}
% ---}

\newcommand{\mydate}[1]{\hrule\vspace{-1mm}\hfill{\small\color{blue}#1}}

% Links {---
\usepackage{hyperref}
\definecolor{darkpowderblue}{rgb}{0.0, 0.2, 0.6}
\hypersetup{colorlinks=true,
    urlcolor=darkpowderblue,
    linkcolor=black,
    citecolor=darkpowderblue,
}
\urlstyle{same}
% ---}


%\usepackage[disable]{todonotes}
\usepackage[draft]{todonotes}

% Comments {---
\usepackage{comment}

% Need this BEFORE specialcomment?
\includecomment{braindump}

% before/after code? Need to set main text back to black...
\specialcomment{braindump}{\small\color{black!50!white}}{\color{black}}

% Comments for everything? Coding, to-do's, science to consider, etc.


%\excludecomment{braindump}
% ---}

\usepackage{caption}
%\renewcommand{\figurename}{}


\begin{document}
%\pagenumbering{roman}
%% Table of contents {---
%\newpage
%\setlength{\parskip}{0.5ex}
%\setcounter{tocdepth}{3}
%\renewcommand\contentsname{}
%\tableofcontents
% ---}
%% Main text {---
%\newpage
%\setlength{\parskip}{2.00ex}
%\newgeometry{left=2.00in,right=0.50in,top=1.0in,bottom=1.0in}
%\setlength{\marginparwidth}{1.50in}
%\setlength{\marginparsep}{0.30in}
%\pagenumbering{arabic}
% ---}
%% Background for margin notes  {---
\begin{comment}
\backgroundsetup{
    color=black,
    opacity=0.1,
    %position={0,-3},
    angle=0,
    scale=4.0,
    hshift=-0.85in,
    contents={
        \begin{tikzpicture}%[x=1.0in, y=1.0in]
            \draw [fill=gray, ultra thick, gray]
            (0,0.1\paperheight) rectangle (1.2,0.9\paperheight);
        \end{tikzpicture}}}
\end{comment}
% ---}




\mydate{10 December 2018}

\paragraph{Recap}
\begin{itemize}
    \item Paper:
        \begin{itemize}
            \item Use mac preview to read paper out loud.\\
                It thinks ``AR11158'' is pronounced ``Arkansas''.
            \item Make sure acronyms are spelled out at first appearance (and
                only then).\\
                This should probably be one of the last things to do before
                submission.
        \end{itemize}
    \item Other:
        \begin{itemize}
            \item Goddard badge\\
                Jack emailed me with a phone number to call about renewal;
                haven't done this yet because I've been in Las Cruces.
        \end{itemize}
\end{itemize}


\paragraph{Agenda}

\begin{itemize}
    \item ``Science questions'' vs. ``Data questions''
    \item Correcting flux for exposure time (see figures on next page).
    \item Subregions - boxed areas in Figure on page 21.
        %\begin{itemize}
        %\end{itemize}
\end{itemize}

\clearpage

\mydate{02 December 2018}

\vspace{1in}

\begin{figure}[htb!]
    \includegraphics[width=\textwidth]{time-3minpower_maps_20181202.pdf}
    \caption*{
        Temporal evolution of 3-minute power (per pixel),
        obtained by summing over the pixel values in each power map
        and dividing by the number of unsaturated pixels.
        The flux was \emph{not} corrected for exposure time.
    }
\end{figure}

\begin{figure}[htb!]
    \includegraphics[width=\textwidth]{time-3minpower_maps.pdf}
    \caption*{
        Same as above, but obtained using flux that \emph{was} corrected
        for exposure time.
    }
\end{figure}




\clearpage


\mydate{01 December 2018}

\vspace{1in}

\begin{figure}[htb!]
    \includegraphics[width=\textwidth]{detrendedLC_20181202.pdf}
    \caption*{
        Raw data light curves overlaid with smoothed, filtered light curves.
        Background = average flux between 00:45 and 01:45.}
\end{figure}

\begin{figure}[htb!]
    \includegraphics[width=\textwidth]{detrended_20181201.pdf}
    \caption*{
        Power spectra from both raw data and detrended data.}
\end{figure}






\clearpage
\mydate{13 August 2018}

%\includegraphics[width=\textwidth]{goes.pdf}
\begin{itemize}
    \item SHINE - spoke with Georgios Chintzoglou about research - suggested
        adding GOES data and updating AIA info (1600\AA{} temperature from Lemen
        et al. 2012 $\rightarrow$ quiet/non-flaring)
    \item Tuesday, August 7 - first day at
        Goddard, badge good through September 30
    \item More background reading on normalizing
        \begin{itemize}
            \item What exactly is the point of doing this?
            \item Normalized FTs of oscillations with higher amplitude causes
                power to be the same in all cases\ldots maybe normalizing helps to
                isolate other contributions to power enhancement/suppression?
        \end{itemize}
    \item Back in Las Cruces at least until Tuesday, August 14
    \item Added GOES data to lightcurves - everything normalized between 0.0
        and 1.0 at the moment
        \begin{itemize}
            \item This took a bit longer than it should have because I
                accidently used November instead of February in tstart/tend.
            \item Milligan et al. 2017 GOES curve looks more
                like 0.5-4.0\AA{} than 1.0-8.0\AA{}. Possible mistake?
        \end{itemize}
    \item Plan for next week:
        \begin{itemize}
            \item update lightcurves with GOES data and complete first draft of
                paper.
            \item Start working with IRIS data? Could first run current
                codes on flare o bserved by IRIS, maybe the one from 10
                September, 2017, or perhaps a smaller one.
        \end{itemize}
\end{itemize}

%\begin{multicols*}{2}
%    \columnbreak
%\end{multicols*}

\end{document}
